\chapter{Foundations: systems, states and state space}

Everything in this subject rests on one idea: that a physical system can be summarised by a short list of
numbers, and that those numbers evolve according to a rule. This chapter builds that idea from nothing.

% ================================================================
\section{What is a system?}

\begin{definition}{System}
A \textbf{system} is anything that takes inputs, has internal behaviour, and produces outputs. Formally it
is a rule mapping an input signal $\uv(t)$ to an output signal $\yv(t)$.
\end{definition}

\begin{intuition}
A robot arm is a system: you command joint targets (input), the motors and links do something (internal
behaviour), and encoders report angles (output). A kettle is a system: power in, temperature out.
\end{intuition}

Systems are classified along axes that matter later:

\begin{itemize}
\item \textbf{Static vs dynamic.} A static system's output depends only on the input \emph{right now}
      ($y = 2u$). A dynamic system's output depends on the \emph{history} of the input --- it has memory.
      A robot arm is dynamic: where it is now depends on every command it has received.
\item \textbf{Linear vs nonlinear.} Linear systems obey superposition: the response to $a u_1 + b u_2$ is
      $a y_1 + b y_2$. Almost nothing real is linear; a great deal is \emph{nearly} linear over a small
      range, which is what makes linearisation (\S\ref{sec:linearisation}) so useful.
\item \textbf{Time-invariant vs time-varying.} Time-invariant means the rule does not change with time:
      the same command today gives the same response tomorrow.
\item \textbf{Continuous-time vs discrete-time.} Continuous time uses derivatives and is how physics is
      written. Discrete time uses difference equations and is how computers run.
\end{itemize}

The abbreviation \textbf{LTI} --- linear time-invariant --- names the class where nearly all the useful
theory lives.

% ================================================================
\section{What is a state?}

This is the central definition of the subject.

\begin{definition}{State}
The \textbf{state} of a system at time $t$ is the smallest set of numbers that, together with all future
inputs, is sufficient to determine all future behaviour. Knowing the state at $t_0$ means the past can be
discarded entirely.
\end{definition}

That phrase --- \emph{the past can be discarded} --- is the whole point. The state is a \emph{summary of
history} that is complete for the purpose of predicting the future.

\begin{intuition}
To predict where a thrown ball goes, you need its position \emph{and} its velocity. Position alone is not
enough --- the same point could be rising or falling. Position and velocity together are enough: gravity
does the rest. So the ball's state is $[\text{position},\ \text{velocity}]$, and how it got there does not
matter.

A useful test: \emph{if I froze the universe and handed you only these numbers, could you resume the
simulation?} If yes, you have a state. If you would need to ask ``but which way was it moving?'', you do
not.
\end{intuition}

\begin{definition}{State vector and state space}
The \textbf{state vector} $\xv \in \R^n$ stacks the state variables into a column. The \textbf{state
space} is the set of all values $\xv$ can take. $n$ is the \textbf{order} of the system.
\end{definition}

\begin{inproject}
The arm has 6 joints. Each joint needs an angle and an angular velocity, exactly like the ball needs
position and velocity. So
\[
\xv = \begin{bmatrix}\qv \\ \dot{\qv}\end{bmatrix} \in \R^{12},
\qquad \qv = [q_1 \ldots q_6]^\top \ \text{(rad)},
\qquad \dot{\qv} \ \text{(rad/s)}.
\]
This is a \textbf{12-state} model. The tracked object has its own state
$[x, y, z, \psi, \dot{x}, \dot{y}, \dot{z}, \dot{\psi}]^\top \in \R^8$: three position components, a yaw
angle, and their rates.
\end{inproject}

\begin{pitfall}
\textbf{``The state is whatever I can measure.''} No --- the state is what is needed to predict the future.
In this project the encoders measure $\qv$ (6 numbers) but the state has 12. The other 6, the joint
velocities, are part of the state and are \emph{never measured}. Recovering them is exactly what the
estimator is for. Conversely, something measurable need not be a state: the camera image has 230{,}400
numbers in it and none of them are states.
\end{pitfall}

% ================================================================
\section{State-space representation}

\begin{definition}{Continuous-time state-space model}
\begin{align}
\dot{\xv}(t) &= \Am\,\xv(t) + \Bm\,\uv(t) \qquad \text{(state equation)}\\
\yv(t) &= \Cm\,\xv(t) + \bm{D}\,\uv(t) \qquad \text{(output equation)}
\end{align}
with $\xv \in \R^n$, $\uv \in \R^m$, $\yv \in \R^p$, and
$\Am \in \R^{n \times n}$, $\Bm \in \R^{n \times m}$, $\Cm \in \R^{p \times n}$, $\bm{D} \in \R^{p \times m}$.
\end{definition}

Each matrix has a plain reading:

\begin{itemize}
\item $\Am$ --- the \textbf{dynamics matrix}. How the system evolves on its own, with no input. Its
      eigenvalues set stability and speed.
\item $\Bm$ --- the \textbf{input matrix}. How the control reaches the states.
\item $\Cm$ --- the \textbf{output matrix}. Which combinations of state the sensors see.
\item $\bm{D}$ --- \textbf{feedthrough}. Input appearing instantly at the output. Usually zero, and zero
      here: a joint position command does not teleport into the encoder reading.
\end{itemize}

\begin{intuition}
Read the first equation aloud as: \emph{the rate of change of the state is a linear combination of where
the state is now, plus a linear combination of what I am commanding.} That is all a state-space model says.
\end{intuition}

\begin{inproject}
With $\xv = [\qv; \dot{\qv}]$ and $\uv$ the joint position targets sent to the inner servo
(\S\ref{sec:servo}, Chapter~2), the model derived in Chapter~9 is
\[
\Am = \begin{bmatrix} \bm{0} & \Id \\[2pt]
      -\Mm^{-1}(\bm{K}_p + \partial \bm{g}/\partial \qv) & -\Mm^{-1}\bm{K}_d \end{bmatrix},
\qquad
\Bm = \begin{bmatrix} \bm{0} \\ \Mm^{-1}\bm{K}_p \end{bmatrix},
\qquad
\Cm = \begin{bmatrix} \Id & \bm{0} \end{bmatrix}.
\]
The top row of $\Am$ says $\dot{\qv} = \dot{\qv}$ --- trivially true, and it is what couples the two halves
of the state together. $\Cm = [\Id\ \ \bm{0}]$ says the encoders measure position and nothing else.
\end{inproject}

\begin{vivabox}
\emph{``Why is the top-right block of $\Am$ the identity?''} Because the derivative of position
\emph{is} velocity. It is definitional, not physical. The physics is all in the bottom row.
\end{vivabox}

% ================================================================
\section{Equilibrium and linearisation}
\label{sec:linearisation}

Real dynamics are nonlinear: $\dot{\xv} = \bm{f}(\xv, \uv)$. The theory above needs matrices. Linearisation
is the bridge.

\begin{definition}{Equilibrium point}
A pair $(\xv_0, \uv_0)$ with $\bm{f}(\xv_0, \uv_0) = \bm{0}$: left alone there, the system stays there.
\end{definition}

\begin{definition}{Linearisation}
Near an equilibrium, expand $\bm{f}$ as a first-order Taylor series. With
$\delta\xv = \xv - \xv_0$ and $\delta\uv = \uv - \uv_0$,
\[
\dot{\delta\xv} \approx
\underbrace{\left.\frac{\partial \bm{f}}{\partial \xv}\right|_{\xv_0,\uv_0}}_{\Am} \delta\xv
+ \underbrace{\left.\frac{\partial \bm{f}}{\partial \uv}\right|_{\xv_0,\uv_0}}_{\Bm} \delta\uv .
\]
$\Am$ and $\Bm$ are \textbf{Jacobian matrices}: matrices of partial derivatives.
\end{definition}

\begin{intuition}
Zoom in far enough on any smooth curve and it looks like a straight line. Linearisation is that, in many
dimensions at once. It is accurate near the point you expanded about and degrades as you move away ---
which is why the project re-computes the linearisation as the arm moves (\emph{gain scheduling},
Chapter~5).
\end{intuition}

\begin{inproject}
Measured accuracy of the linearisation, from \code{scripts/model\_analysis.py}: relative RMS error against
the full nonlinear simulator is \textbf{0.12\,\%} over a 0.02\,rad step and \textbf{2.1\,\%} over a
0.3\,rad step. The error grows with step size exactly as the theory predicts. This is also the check that
caught the servo-damping defect --- see Chapter~9.
\end{inproject}

\begin{pitfall}
\textbf{``Linearisation fails when the system is nonlinear.''} It fails when the system is nonlinear
\emph{over the range you are using it}. The dangerous case is not curvature but a \emph{hard} nonlinearity
like saturation: past the limit, the slope is simply wrong, and no amount of ``staying close'' helps. That
is precisely the defect in Chapter~9 --- torque saturation, not curvature, broke the model.
\end{pitfall}

% ================================================================
\section{Discretisation}

Physics is continuous; controllers run on computers at a fixed rate. Discretisation converts one to the
other.

\begin{definition}{Zero-order hold (ZOH)}
The input is held constant between sampling instants: $\uv(t) = \uv_k$ for $t \in [k\,dt, (k{+}1)dt)$.
Under ZOH the exact discrete equivalent of $\dot{\xv} = \Am\xv + \Bm\uv$ is
\[
\xv_{k+1} = \Am_d \xv_k + \Bm_d \uv_k, \qquad
\Am_d = e^{\Am\,dt}, \qquad
\Bm_d = \left(\int_0^{dt} e^{\Am\tau}d\tau\right)\Bm .
\]
\end{definition}

Both are computed in one step by exponentiating a block matrix:
\[
\exp\!\left(\begin{bmatrix} \Am & \Bm \\ \bm{0} & \bm{0}\end{bmatrix} dt\right)
= \begin{bmatrix} \Am_d & \Bm_d \\ \bm{0} & \Id \end{bmatrix}.
\]

\begin{definition}{Matrix exponential}
$e^{\Am} = \sum_{k=0}^{\infty} \Am^k / k!$ --- the matrix generalisation of the scalar exponential. It
solves $\dot{\xv} = \Am\xv$ exactly: $\xv(t) = e^{\Am t}\xv(0)$.
\end{definition}

A cruder alternative is the \textbf{forward Euler} approximation $\Am_d \approx \Id + \Am\,dt$, which is
the first two terms of that series. It is only accurate when $|\lambda|\,dt \ll 1$ for every eigenvalue.

\begin{inproject}
Control period $dt = 10$\,ms (100\,Hz); the physics runs at 500\,Hz, so there are 5 physics substeps per
control step. The fastest closed-loop pole is $-161$\,s$^{-1}$, giving $|\lambda|\,dt \approx 1.6$ ---
\emph{not} small. Euler would be badly wrong here, so the project uses the matrix exponential
(\code{scipy.linalg.expm}) in \code{core/kinematics.py:discrete\_model}.
\end{inproject}

\begin{vivabox}
\emph{``Your poles are at $-161$ and you sample at 100\,Hz. Is that fast enough?''} With ZOH, yes, but it
is marginal: $|\lambda| dt \approx 1.6$. It is one of the three things that bound how far the controller
gain can be pushed (Chapter~5). And it is why Euler discretisation was not an option.
\end{vivabox}

% ================================================================
\section{Stability and poles}

\begin{definition}{Eigenvalue, eigenvector}
$\lambda$ and $\bm{v} \neq \bm{0}$ satisfying $\Am\bm{v} = \lambda\bm{v}$. The matrix acts on $\bm{v}$ as
pure scaling. Eigenvalues of a system matrix are called \textbf{poles}.
\end{definition}

\begin{definition}{Asymptotic stability}
Continuous time: stable iff every eigenvalue of $\Am$ has $\mathrm{Re}(\lambda) < 0$.\\
Discrete time: stable iff every eigenvalue of $\Am_d$ has $|\lambda| < 1$ (inside the unit circle).
\end{definition}

\begin{intuition}
Decompose the state along eigenvectors; each component then evolves independently as $e^{\lambda t}$.
Negative real part means decay. A positive real part means at least one direction grows without bound ---
divergence. The imaginary part sets oscillation frequency, and the magnitude of the real part sets how
fast that mode settles: time constant $\tau = 1/|\mathrm{Re}(\lambda)|$.
\end{intuition}

\begin{inproject}
Closed-loop poles span $-161$ to $-26$\,s$^{-1}$ across the tested configurations: all stable, with time
constants between 6 and 38\,ms. The spread is why the slowest joints dominate settling and why the
reference filter is tuned around $\omega = 5$\,rad/s, far below the fastest pole.
\end{inproject}

% ================================================================
\section{Controllability}

\begin{definition}{Controllability}
A system is \textbf{controllable} if for any initial state and any target state there exists an input that
drives one to the other in finite time. The test is the rank of the \textbf{controllability matrix}
\[
\mathcal{C} = \begin{bmatrix}\Bm & \Am\Bm & \Am^2\Bm & \cdots & \Am^{n-1}\Bm\end{bmatrix},
\]
which must be $n$ (full row rank).
\end{definition}

\begin{intuition}
Can my actuators reach every corner of the state space? Each extra block $\Am^k\Bm$ asks: having pushed,
and let the dynamics propagate that push, which new directions have opened up? If after $n{-}1$ rounds some
direction is still unreachable, no input will ever reach it.

A car with a stuck steering wheel is uncontrollable in lateral position: the engine cannot produce sideways
motion no matter how it is used.
\end{intuition}

\begin{definition}{Minimum singular value as a conditioning measure}
Rank is binary and numerically fragile. $\sigma_{\min}(\mathcal{C})$ measures \emph{how far} the matrix is
from losing rank --- how strongly, not merely whether, the system is controllable.
\end{definition}

\begin{inproject}
Rank $12/12$ at every configuration tested, with $\sigma_{\min} = 0.320$. Each joint has its own actuator,
so this is expected --- but it is tested rather than assumed, in
\code{scripts/model\_analysis.py}.
\end{inproject}

% ================================================================
\section{Observability}

The dual question, and the one that matters most in this project.

\begin{definition}{Observability}
A system is \textbf{observable} if the initial state can be determined uniquely from a finite record of
outputs and inputs. The test is the rank of the \textbf{observability matrix}
\[
\mathcal{O} = \begin{bmatrix}\Cm \\ \Cm\Am \\ \Cm\Am^2 \\ \vdots \\ \Cm\Am^{n-1}\end{bmatrix},
\]
which must be $n$ (full column rank).
\end{definition}

\begin{intuition}
Can I work out everything about the system from what I can see? Not instantaneously --- \emph{over time}.
Each block row asks what an extra derivative of the measurement reveals. You cannot see a ball's velocity
in a single photograph, but two photographs give it to you. The observability matrix formalises exactly
that argument.
\end{intuition}

\begin{inproject}
This is the crux of the whole project. $\Cm = [\Id\ \ \bm{0}]$: encoders measure 6 of 12 states, so half
the state is invisible. Yet
\[
\Cm\Am = \begin{bmatrix}\Id & \bm{0}\end{bmatrix}
\begin{bmatrix}\bm{0} & \Id \\ \ast & \ast\end{bmatrix} = \begin{bmatrix}\bm{0} & \Id\end{bmatrix},
\]
so the second block row of $\mathcal{O}$ sees velocity directly --- because the derivative of a measured
position \emph{is} the velocity. Rank is $12/12$: the velocities are observable, and so an estimator can
reconstruct them.

Measured results:
\begin{center}
\begin{tabular}{@{}lll@{}}
\toprule
Sensor set & Rank & $\sigma_{\min}$ \\
\midrule
All 6 encoders & $12/12$ & $4.5 \times 10^{-3}$ \\
Single encoder (shoulder pan only) & $8$--$9 / 12$ & --- \\
\bottomrule
\end{tabular}
\end{center}
Full rank, but $\sigma_{\min}$ is small: the system is \textbf{weakly observable}. That weakness has a
practical consequence --- the velocity estimate is the part that degrades first when the filter is tuned
badly (Chapter~4).
\end{inproject}

\begin{vivabox}
\emph{``Which states are observable from which sensors, and what happens when the object leaves view?''}
A named exam question. Answer in two halves. \textbf{Joint states:} rank $12/12$ from encoders alone,
$\sigma_{\min} = 4.5\times10^{-3}$, so observable but weakly --- and that is exactly why velocity degrades
first under a mis-set $\Qm$. With a single encoder the rank falls to $8$--$9$. \textbf{Object states:}
position and yaw are measured when in view; velocity never is, and is reconstructed. When the camera is
lost the filter stops correcting and \emph{coasts} on its process model --- see Chapter~4.
\end{vivabox}

\begin{definition}{Duality}
Controllability of $(\Am, \Bm)$ is mathematically identical to observability of $(\Am^\top, \Bm^\top)$.
Every theorem about one has a mirror image about the other.
\end{definition}

\begin{intuition}
This is why LQR (optimal control) and the Kalman filter (optimal estimation) have the same algebraic
form --- both solve a Riccati equation. They are the same problem viewed from opposite ends: one pushes
the state where you want it, the other infers where it is.
\end{intuition}
